\section{Introducción: tres hilos de lectura de K2.5}
\begin{figure}[H]
\centering
\includegraphics[width=0.78\textwidth]{cover.jpg}
\caption{La imagen de ``Visual Agency Intelligence'' de K2.5 en el adelanto de lanzamiento, que enfatiza el punto de partida de la capacidad de agente visual. \protect\footnotemark}
\end{figure}
\footnotetext{Intervalo de tiempo en el video: 00:00:05--00:00:13, la ruta técnica de K2.5 que señala el docente en la apertura.}
Moonshot AI presentó oficialmente K2.5 el 25/26 de enero de 2026, y el docente dijo desde la apertura: \textit{``Dive into Kimi K2.5 这篇庞杂的文章''}, y ubicó todo este curso como un recorrido integral desde los datos y el algoritmo hasta la infraestructura. También se menciona en esta misma línea de tiempo la versión de Gemini 3 Flash agent vision, publicada dos días después, con el fin de mostrar que la multimodalidad nativa y el agente visual no son eventos aislados.

\begin{knowledgebox}{Tres hilos de lectura}
\begin{itemize}
    \item \textbf{Data}: indagar cómo se define, se muestrea y se sintetiza cada nuevo dominio de datos; un nuevo domain implica una nueva capacidad.
    \item \textbf{Algorithm}: ¿quién evalúa la respuesta del modelo? Cómo GRM y el entrenamiento en varias etapas colaboran para estabilizar las tareas open-ended.
    \item \textbf{Infra}: cómo Critical Steps, el agent swarm en paralelo y el flujo de evaluación gestionan juntos el compute y la latency.
\end{itemize}
\end{knowledgebox}

\subsection{Resumen de la sección}
La introducción a K2.5 parte de ``Read the tech report'', deja claro que usaremos las tres perspectivas de datos, algoritmo y arquitectura para seguir la ruta de implementación de Visual Agency Intelligence, y toma este guion como la versión anotada del informe técnico.
